\documentclass[conference]{IEEEtran}
\usepackage{cite,amsmath,graphicx,booktabs,array,url,balance}
\usepackage[T1]{fontenc}
\hyphenation{multi-sensor orches-trator}
\begin{document}
\title{Agentic AI for Intelligent Border Surveillance: A Multi-Agent Framework for Automatic Threat Detection, Risk Prioritization, and Controlled Human-in-the-Loop Response}
\author{\IEEEauthorblockN{[Your Name]}
\IEEEauthorblockA{[Department], [Institution], [City, Country]\\
Email: [your.email@institution.edu]}}
\maketitle

\begin{abstract}
Border areas are long, remote, and exposed to infiltration, vehicle intrusion, hostile drones, and fence tampering, yet most monitoring still depends on operators watching many sensor feeds continuously. This paper proposes an agentic framework in which specialised software agents perceive, fuse, reason about, and prioritise events, while a central orchestrator controls what each agent may do and keeps a human operator in charge of consequential decisions. The framework detects unauthorised crossings, loitering, group movement, vehicle intrusion, drones, abandoned objects, and fence or gate breaches; assigns each event a risk level (Low, Medium, High, Critical); and issues real-time alerts ranked by priority. We evaluate the pipeline in a 60-day synthetic multi-sensor simulation with a camera-only baseline and a multi-sensor OR rule. On the 30-day test period, the proposed fusion reduced false alarms from 33.0 to 1.6 per day and operator interruptions from 38.3 to 6.3 per day relative to the OR rule, while retaining 91.8\% threat recall and 98.1\% risk-level accuracy. These results come from simulated data under stated assumptions and show the behaviour of the design; they are not field measurements, and field validation is required.
\end{abstract}
\begin{IEEEkeywords}
border surveillance, agentic AI, multi-agent systems, sensor fusion, anomaly detection, risk assessment, human-in-the-loop
\end{IEEEkeywords}

\section{Introduction}
Securing a border means watching hundreds or thousands of kilometres of rough terrain, in darkness, fog, and dust, with limited staff. Fixed cameras, radar, and ground sensors generate far more alarms than people can verify. Animals, wind-blown vegetation, and authorised patrols trigger the same sensors as real intruders, so operators lose time on false alarms and risk missing the real events. Continuous manual monitoring is also tiring, and fatigue lengthens response time.

Deep learning has improved object detection \cite{redmon2016,ren2015} and tracking \cite{bewley2016,wojke2017}, but a detector alone does not give a security team what it needs. Staff need to know \emph{what} happened, \emph{how serious} it is, \emph{what to do first}, and \emph{why the system believes it}. This calls for a system that combines several sensors, reasons about behaviour and context, ranks events, and acts only inside clearly defined limits. Agent-based designs are a natural fit \cite{wooldridge1995,wang2024}: each agent has a narrow role, and an orchestrator coordinates them.

This paper makes four contributions:
\begin{enumerate}
\item A multi-agent architecture covering nine operational objectives: automatic threat detection, infiltration and crossing detection, suspicious movement/loitering/group behaviour, vehicle intrusion, drone and abandoned-object detection, fence/gate/zone monitoring, real-time alerts, risk levels, and reduced operator workload with faster response.
\item A control design (orchestrator, policy layer, human approval, audit log, fail-safe fallback) that bounds what the agents can do.
\item A transparent risk-scoring and escalation scheme that turns detections into prioritised alerts.
\item A reproducible simulation study comparing the framework with two baselines, with explicit limitations.
\end{enumerate}

\section{Related Work}
\textbf{Detection and tracking.} Single-stage and two-stage detectors \cite{redmon2016,ren2015} and tracking-by-detection \cite{bewley2016,wojke2017} are standard building blocks for people and vehicle monitoring. Drone-view benchmarks \cite{zhu2022} show that small, fast objects remain difficult.

\textbf{Sensor fusion.} Combining radar, thermal, visual, and vibration sensors lowers false alarms because their failure modes differ \cite{hall1997}. Our fusion agent follows this principle but adds behavioural and contextual evidence.

\textbf{Anomaly and behaviour analysis.} Loitering, unusual paths, and abnormal dwell times are classical anomaly-detection targets \cite{chandola2009}. Rare events and scarce labels make supervised learning alone inadequate.

\textbf{Agentic systems.} Agents that perceive, reason, and act through tools \cite{wooldridge1995,yao2023,wang2024} offer modularity but raise control questions: which actions are permitted, who approves them, and how decisions are audited. Guidance for human--AI interaction stresses clear scope, explanation, and correction \cite{amershi2019}. We address these questions explicitly in Section~IV-C.

\section{Requirements}
Table~\ref{tab:req} maps each stated objective to the component that satisfies it and the metric used to judge it.
\begin{table}[!t]
\caption{Objectives, Responsible Components, and Metrics}\label{tab:req}
\centering\footnotesize
\begin{tabular}{p{2.55cm}p{2.7cm}p{2.2cm}}
\toprule
\textbf{Objective} & \textbf{Component} & \textbf{Metric}\\
\midrule
Automatic detection; crossings & Vision, Fusion & Recall, AUC\\
Loitering, group movement & Behaviour agent & Per-class recall\\
Vehicle intrusion & Vehicle, Fusion & Per-class recall\\
Drones, abandoned objects & Drone, Vision & Per-class recall\\
Fence/gate/zone intrusion & Perimeter agent & Per-class recall\\
Real-time alerts & Alert agent & Mean queue wait\\
Risk levels & Risk agent & Risk accuracy\\
Less human workload & Orchestrator & Interruptions/day\\
Fewer false alarms & Fusion & False alarms/day\\
\bottomrule
\end{tabular}
\end{table}

\section{Proposed Agentic Framework}
\subsection{Architecture}
Fig.~\ref{fig:arch} shows the five-stage pipeline. Heterogeneous sensors (visible and thermal cameras, radar, fence/fibre sensing, RF/acoustic drone detection) feed \emph{perception agents} running at the edge, which detect and track targets close to the sensor so that only compact events cross the network \cite{shi2016}. \emph{Reasoning agents} then fuse the evidence, analyse behaviour, and score risk. The \emph{orchestrator} decides what happens next, and an \emph{alert agent} delivers prioritised messages to a human operator. Operator confirmations and corrections are logged and used to retrain the models.
\begin{figure*}[!t]
\centering\includegraphics[width=0.92\textwidth]{figs/fig1_arch.pdf}
\caption{Proposed multi-agent architecture. Perception agents run at the edge; reasoning agents fuse evidence, analyse behaviour, and score risk; the orchestrator applies policy before alerts reach the operator. The dashed arrow is the retraining feedback loop.}\label{fig:arch}
\end{figure*}

\subsection{Agent Roles}
Table~\ref{tab:agents} lists each agent's inputs, outputs, and method. The framework does not fix the model inside an agent; detectors, trackers, or rule modules can be replaced without changing the interfaces.
\begin{table}[!t]
\caption{Agent Roles}\label{tab:agents}
\centering\footnotesize
\begin{tabular}{p{1.35cm}p{2.9cm}p{3.0cm}}
\toprule
\textbf{Agent} & \textbf{Task} & \textbf{Output}\\
\midrule
Vision & Detect/track people, groups, objects in visible and thermal video & Tracks, class, confidence\\
Vehicle & Detect vehicles; speed and path; match authorised patrols & Vehicle events\\
Drone & Radar and RF/acoustic drone detection & Drone events, bearing\\
Perimeter & Fence vibration, gate state, restricted-zone entry & Breach events\\
Fusion & Combine sensors with weather, time, patrol schedule & Threat probability per class\\
Behaviour & Loitering, dwell, group size, heading & Behaviour features\\
Risk & Severity and proximity scoring & Risk level L/M/H/C\\
Alert & Console, digest, radio/SMS & Prioritised alerts\\
\bottomrule
\end{tabular}
\end{table}

\subsection{Control of the Agents}
The system is designed so that the agents advise and the humans decide. Six mechanisms enforce this.
\begin{enumerate}
\item \textbf{Orchestrator and typed messages.} Agents communicate only through the orchestrator using a fixed event schema (sensor, time, location, class, confidence). No agent calls another directly, so every decision passes one control point.
\item \textbf{Tool allow-list.} Each agent can use only the tools its role needs. Perception agents can read sensors and cue a camera; none can operate barriers, weapons, or vehicles. Launching a drone for inspection requires operator approval.
\item \textbf{Escalation policy.} The risk level determines delivery (Table~\ref{tab:esc}). Critical and High events require human acknowledgement before any response is dispatched.
\item \textbf{Confidence and context checks.} Low-confidence events are held for corroboration by a second sensor within a time window; events matching the authorised patrol schedule are down-ranked, not deleted.
\item \textbf{Audit and feedback.} Every event, score, and operator action is logged. Operator corrections become labelled training data.
\item \textbf{Fail-safe.} If an agent or link fails, the orchestrator reverts to raw sensor alarms for that zone and notifies the operator, so a software fault cannot silence a sensor.
\end{enumerate}
\begin{table}[!t]
\caption{Escalation Policy by Risk Level}\label{tab:esc}
\centering\footnotesize
\begin{tabular}{p{1.1cm}p{3.0cm}p{3.1cm}}
\toprule
\textbf{Level} & \textbf{Delivery} & \textbf{Human role}\\
\midrule
Critical & Immediate alert, audio, map zoom, evidence clip & Must acknowledge and approve response\\
High & Real-time alert, top of queue & Must acknowledge\\
Medium & Real-time queue & Review in order\\
Low & Periodic digest & Review, relabel\\
\bottomrule
\end{tabular}
\end{table}

\subsection{Risk Scoring}
For an event with feature vector $\mathbf{x}$, the fusion agent outputs class probabilities $P(c\mid\mathbf{x})$ over threat classes $c\in\mathcal{T}$ and benign classes. The threat score is $S=\sum_{c\in\mathcal{T}}P(c\mid\mathbf{x})$, and an alert is raised when $S\ge\tau$. The risk agent computes
\begin{equation}
R=\bar{\sigma}\,\big(1+0.3\,\mathbb{1}[g\ge4]+0.4\,\mathbb{1}[d<100]+0.2\,\mathbb{1}[n]\big),
\end{equation}
where $\bar{\sigma}=\sum_{c\in\mathcal{T}}P(c\mid\mathbf{x})\sigma_c\big/\sum_{c\in\mathcal{T}}P(c\mid\mathbf{x})$ is the probability-weighted class severity (with $\sigma_c\in[0.3,0.9]$ per class), $g$ the group size, $d$ the distance to the border in metres, and $n$ indicates night. $R$ is mapped to Low ($<0.45$), Medium ($<0.70$), High ($<1.0$), or Critical. The weights are design parameters that an agency should set from its own doctrine; they are not learned.

Fig.~\ref{fig:scene} illustrates the output on a border segment: tracks are coloured by risk, and a benign animal track is kept at Low rather than interrupting the operator.
\begin{figure}[!t]
\centering\includegraphics[width=0.78\columnwidth]{figs/fig2_scene.pdf}
\caption{Schematic top-down view (illustrative, not real data). Sensor towers T1--T3 with fields of view, radar, a gate, and a restricted zone. Tracks are coloured by assigned risk level.}\label{fig:scene}
\end{figure}

\section{Experimental Setup}
\textbf{Why simulation.} Operational border data are restricted and rarely released. We therefore built a synthetic event generator to test the logic of the pipeline. All numbers below are produced by this simulation.

\textbf{Events.} The generator produces 60 days of events. Seven threat classes (person crossing, group, loitering, vehicle intrusion, drone, abandoned object, fence breach) arrive at 6.0 events/day in total. Four benign classes (animal, vegetation motion, authorised patrol, weather clutter) arrive at 65/day. Threats are weighted toward night (60\%), and 20\% of events occur in bad weather.

\textbf{Sensors.} Each event yields a noisy score in $[0,1]$ from five sensors (camera, thermal, radar, fence, RF). Class-dependent means are assumed (e.g., the camera responds to vegetation, radar to animals and vehicles, and RF almost only to drones). The camera mean is reduced by 55\% at night and 25\% in bad weather. Behavioural features (speed, dwell time, group size, heading, distance to border) and an authorised-patrol flag are also generated per class.

\textbf{Methods.} (i) \emph{Camera-only}: alarm if the camera score exceeds 0.5. (ii) \emph{Multi-sensor OR}: alarm if any sensor exceeds 0.6. (iii) \emph{Proposed}: a gradient-boosted classifier \cite{friedman2001} trained on days 1--30 over all sensor, behaviour, and context features, evaluated on days 31--60 (171 threat and 1,971 benign events). The threshold $\tau$ was set on the training data by five-fold cross-validation to reach at least 92\% recall.

\textbf{Operator model.} One operator verifies each real-time alert in an exponentially distributed time with mean 5 min. The baselines serve alerts first-come-first-served. The proposed system serves Medium and above in priority order and sends Low events to a digest.

\textbf{Metrics.} ROC-AUC, recall, precision, false alarms/day, operator interruptions/day, recall of High/Critical events, mean queue wait of High/Critical events, and risk-level accuracy.

\section{Results}
\textbf{Detection.} Fig.~\ref{fig:roc}(a) shows that the camera-only alarm is weak (AUC 0.576) because it misses night and drone events and fires on vegetation. The OR rule reaches AUC 0.858 but cannot separate threats from animals and patrols. The fusion agent reaches AUC 0.990. At the chosen operating point (Table~\ref{tab:res}), it detects 91.8\% of threats with 76.2\% precision; the OR rule has similar recall (93.6\%) but 13.9\% precision.

\begin{figure*}[!t]
\centering\includegraphics[width=0.9\textwidth]{figs/fig3_roc.pdf}
\caption{(a) ROC curves on the 30-day test set. (b) Recall for each threat class at each method's operating point.}\label{fig:roc}
\end{figure*}

\textbf{Per-class behaviour.} In Fig.~\ref{fig:roc}(b) the camera-only baseline fails on drones (9\% recall) and fence breaches (19\%). The fusion agent's recall is 100\% for drones, abandoned objects, and fence breaches, 92--96\% for persons, loitering, and vehicles, but only 60\% for groups. Groups are confused with persons and animals when the camera is degraded at night, so this class is the weakest point of the current design.

\begin{table}[!t]
\caption{Test-Period Results (30 Days, Simulated Data)}\label{tab:res}
\centering\footnotesize
\setlength{\tabcolsep}{3.2pt}
\begin{tabular}{lccc}
\toprule
\textbf{Metric} & \textbf{Camera} & \textbf{OR rule} & \textbf{Proposed}\\
\midrule
ROC-AUC & 0.576 & 0.858 & \textbf{0.990}\\
Threat recall (\%) & 42.1 & \textbf{93.6} & 91.8\\
Precision (\%) & 11.7 & 13.9 & \textbf{76.2}\\
False alarms/day & 18.2 & 33.0 & \textbf{1.6}\\
Interruptions/day & 20.6 & 38.3 & \textbf{6.3}\\
High/Crit. recall (\%) & 38.5 & \textbf{93.4} & 91.0\\
Mean wait, High/Crit. (s) & 29 & 39 & $<$1\\
Risk-level acc. (\%) & -- & -- & 98.1\\
\bottomrule
\end{tabular}
\end{table}

\begin{figure*}[!t]
\centering\includegraphics[width=0.9\textwidth]{figs/fig4_ops.pdf}
\caption{(a) False alarms per day. (b) Alerts that interrupt the operator per day. (c) Risk-level confusion matrix for correctly detected threats (L, M, H, C = Low, Medium, High, Critical).}\label{fig:ops}
\end{figure*}

\textbf{Workload.} Relative to the OR rule, the proposed system cuts false alarms by 95\% (33.0 to 1.6 per day) and operator interruptions by 84\% (38.3 to 6.3 per day); relative to camera-only, interruptions fall by 69\% (Fig.~\ref{fig:ops}a,b). A further 0.6 Low-risk events per day go to the digest rather than the live queue.

\textbf{Risk levels.} Of the detected threats, 98.1\% received the correct level and 99.4\% were within one level (Fig.~\ref{fig:ops}c). This figure is high partly because the generator uses the same features for the true and estimated risk, so it tests consistency rather than expert agreement.

\textbf{Response time.} With a single operator at 5 min per verification, queue waits are short in all cases (29 s, 39 s, and under 1 s), because the alert rates are modest. The benefit of prioritisation would be larger when alerts arrive in bursts or operators are scarce; we did not simulate those conditions.

\textbf{Trade-off.} The OR rule keeps slightly higher High/Critical recall (93.4\% against 91.0\%) at a cost of about 20 times more false alarms. About 9\% of serious events were missed by the fusion agent in this setting, so a deployment should tune $\tau$ toward recall and accept more alerts if the cost of a miss is high.

\section{Discussion and Limitations}
\textbf{Validity.} The simulated test data come from the same generative model as the training data, so the fusion agent has an advantage that real deployment will not offer. Sensor responses, event rates, and operator times are assumptions. Real terrain will bring domain shift, sensor drift, and adversarial behaviour such as camouflage, thermal masking, spoofing, and jamming. The results should be read as showing that the design works as intended under idealised conditions, not as expected field accuracy. The next step is to replace the generator with recorded sensor data from a test range, using real detectors in the perception agents.

\textbf{Not evaluated.} We did not evaluate language-model agents, vision models on real imagery, communication outages, or cybersecurity. Any language-model component, for example an operator-facing summary, should operate only on structured events and stay under the same allow-list and policy checks.

\textbf{Ethics.} Border systems affect civilians, migrants, and local communities. Misclassification can cause harm, so the framework withholds any autonomous use of force, logs every decision, and makes risk weights visible and changeable. Data retention, privacy, bias testing across conditions, and independent audit are needed before deployment, and the same sensors can support search and rescue for people in distress.

\section{Conclusion}
We presented a multi-agent framework in which perception, fusion, behaviour, risk, and alert agents are coordinated by a policy-bound orchestrator with human approval for serious events. In simulation, it removed most false alarms and cut operator interruptions by 84\% relative to a multi-sensor OR rule while keeping above 90\% threat recall and assigning risk levels consistently. Group movement under degraded night vision was the weakest class. Future work will test the framework on recorded multi-sensor data, add temporal models for behaviour, study burst conditions and multi-operator queues, and evaluate robustness against adversarial attacks.

\begin{thebibliography}{99}\small
\bibitem{redmon2016} J. Redmon, S. Divvala, R. Girshick, and A. Farhadi, ``You only look once: Unified, real-time object detection,'' in \emph{Proc. IEEE CVPR}, 2016, pp. 779--788.
\bibitem{ren2015} S. Ren, K. He, R. Girshick, and J. Sun, ``Faster R-CNN: Towards real-time object detection with region proposal networks,'' in \emph{Proc. NeurIPS}, 2015.
\bibitem{bewley2016} A. Bewley, Z. Ge, L. Ott, F. Ramos, and B. Upcroft, ``Simple online and realtime tracking,'' in \emph{Proc. IEEE ICIP}, 2016, pp. 3464--3468.
\bibitem{wojke2017} N. Wojke, A. Bewley, and D. Paulus, ``Simple online and realtime tracking with a deep association metric,'' in \emph{Proc. IEEE ICIP}, 2017, pp. 3645--3649.
\bibitem{wooldridge1995} M. Wooldridge and N. R. Jennings, ``Intelligent agents: Theory and practice,'' \emph{Knowledge Eng. Rev.}, vol. 10, no. 2, pp. 115--152, 1995.
\bibitem{wang2024} L. Wang \emph{et al.}, ``A survey on large language model based autonomous agents,'' \emph{Front. Comput. Sci.}, vol. 18, no. 6, 2024.
\bibitem{yao2023} S. Yao \emph{et al.}, ``ReAct: Synergizing reasoning and acting in language models,'' in \emph{Proc. ICLR}, 2023.
\bibitem{hall1997} D. L. Hall and J. Llinas, ``An introduction to multisensor data fusion,'' \emph{Proc. IEEE}, vol. 85, no. 1, pp. 6--23, Jan. 1997.
\bibitem{chandola2009} V. Chandola, A. Banerjee, and V. Kumar, ``Anomaly detection: A survey,'' \emph{ACM Comput. Surv.}, vol. 41, no. 3, pp. 1--58, 2009.
\bibitem{shi2016} W. Shi, J. Cao, Q. Zhang, Y. Li, and L. Xu, ``Edge computing: Vision and challenges,'' \emph{IEEE Internet Things J.}, vol. 3, no. 5, pp. 637--646, Oct. 2016.
\bibitem{zhu2022} P. Zhu \emph{et al.}, ``Detection and tracking meet drones challenge,'' \emph{IEEE Trans. Pattern Anal. Mach. Intell.}, vol. 44, no. 11, 2022.
\bibitem{friedman2001} J. H. Friedman, ``Greedy function approximation: A gradient boosting machine,'' \emph{Ann. Statist.}, vol. 29, no. 5, pp. 1189--1232, 2001.
\bibitem{amershi2019} S. Amershi \emph{et al.}, ``Guidelines for human-AI interaction,'' in \emph{Proc. ACM CHI}, 2019.
\end{thebibliography}
\end{document}
